% Pursuing Stacks (cotes 134-2 to 134-8): archive cote 134-2, site page 17 (PDF page 18).
% letter to Quillen. Draft compared: ps1.tex:97-179.
% AI-transcribed by Claude Opus 5.5 (claude-opus-5-5, adaptive thinking, effort high), Message Batches API, batch pilot (msgbatch_0189Nhf1xB615YWnztubtLfc), 2026-11-12; not yet checked by eye.
\scanpage{134-2/17}
\mspage{2)}
\margin{3}

had been done and the homotopy category expressed in terms of $\infty$-groupoids\unsure{.}
But finally it appears this is not so,  they have been working throughout
with a \unsure{nction|notion} of $\infty$-groupoid too restrictive for the purposes I had in
mind (probably because they insist I guess on strict associativity of
compositions,  rather than associativity up to a (given) isomorphism\inserted{(or rather, homotopy)}) -
to the effect that the simply connected homotopy types they obtain are
merely products of Eilenberg-MacLane spaces,  too bad !  They do not seem
to have realized yet that this makes their set-up wholly unadequate \deleted{for} \hand{as}
a sweeping foundational set-up for homotopy.  This brings to the fore
again to work out the suitable definitions for n-groupoids -  if this is
not done yet anywhere.  I spent the afternoon today trying to figure out
a reasonable definition,  to get a feeling at least of where the difficul-%
ties are, if any.  I am guided mainly of course by the topological inter-%
pretation.  It will be short enough to say how far I got.  The main part
of the structure it seems is expressed by the sets $F_i$ $(i\unsure{\in}\underline{N})$ of i-objects,
the source, target and identity maps
\[ s^i_1, t^i_1 : F_i \longrightarrow F_{i-1} \qquad (i \geq 1) \]
\[ k^i_1 : F_i \longrightarrow F_{i+1} \qquad (i \in \underline{N}) \quad \hand{, notation $id_a = k^i_1(a)$,} \]
and the symmetry map (passage to the inverse) \hand{$\check{u} = inv_i(u)$}
\[ inv_i : \unsure{(}F_i \longrightarrow F_i \qquad (i \geq 1) \;, \]
satisfying some obvious relations :  $k^i_1$ is right inverse to the \xout{target}
source and target maps $s^{i+1}_1$ , $t^{i+1}_1$ ,  $inv_i$  is an involution and ``\unsure{exehan|exchan}-%
ges'' source and target,  and moreover for  $i \geq 2$
\[ s^{i-1}_1 s^i_1 = s^{i-1}_1 t^i_1 \quad (\overset{def}{=} s^i_2 : F_i \longrightarrow F_{i-2}) \]
\[ t^{i-1}_1 s^i_1 = t^{i-1}_1 t^i_1 \quad (\overset{def}{=} t^i_2 : F_i \longrightarrow F_{i-2}) \quad ; \]
thus the composition of the source and target maps yields,  for  $0 \leqslant j \leqslant i$ ,
just \uline{two} maps
\[ s^i_{\hand{\ell}} , t^i_{\hand{\ell}} : F_i \rightarrow F_{i-\hand{\ell}} \hand{= F_j} \; ( \hand{\ell} = i-j ) . \]
The next basic structure is the composition structure,  where the usual
composition of arrows,  more specifically of i-objects (i$\geq$1)  $v \circ u$
(defined when $t_1(u)=s_1(v)$)  must be supplemented by the Godement-type
operations $\hand{\mu} * \hand{\lambda}$ when $\hand{\mu}$ and $\hand{\lambda}$ are``arrows between arrows''\inserted{, etc}.  Following
this line of thought,  one gets the composition maps
\[ (u,v) \longmapsto ( v \overset{\hand{i}}{\underset{\hand{\ell}}{*}} u \quad : \quad (F_i, \unsure{s}^i_{\hand{\ell}})_{u,v} \times_{F_{i-\hand{\ell}}} (F_i, s^i_{\hand{\ell}}) \longrightarrow F_i \; , \]
\margin{$\overset{i}{\underset{\ell}{*}}$ / $\underset{\ell}{*}$}
the \deleted{\illegible{1}} \inserted{$\overset{i}{\underset{\ell}{*}}$-composition of} \xout{operation} i-objects ( $1 \leqslant \hand{\ell} \leqslant i$)  being defined when the
$\hand{\ell}$-target of u is equal to the $\hand{\ell}$-source of v,  and then we have
\[ \deleted{\illegible{1}} \; s^i_1(v \underset{\hand{\ell}}{*} u) \; = \; s^i_1(v) \underset{\hand{\ell-1}}{*} s^i_1(u) \]
\[ \hand{t^i_1(v \underset{\ell}{*} u) = t^i_1(v) \underset{\ell-1}{*} s^i_1(u)} \quad \hand{\Big\}\ \ell \geq 2 \text{ i.e. } \ell-1 \geq 1} \]

\margin{17}
